\subsubsection{The persistent return of the five regions}

Let
\[
\mathcal F_\pi=\Psi^{-1}(010211)\subset\mathcal U_6
\]
be the microfibre relative to the catalogue. Its five elements are recalled
below, with each word separated into an initial block and a terminal block:
\[
\begin{aligned}
 &(501,614,498\mid169,272,272),\\
 &(810,923,870\mid169,272,272),\\
 &(810,983,810\mid169,272,272),\\
 &(870,923,810\mid169,272,272),\\
 &(870,923,810\mid418,674,521).
\end{aligned}
\]
Define the return projection
\[
p_{\rm ret}:\mathbb Z^6\longrightarrow\mathbb Z^3,
\qquad
p_{\rm ret}(u_1,\ldots,u_6)=(u_4,u_5,u_6).
\]

\begin{proposition}[Persistent terminal block]
\label{p12-83-c13-s3:prop:bloque-terminal-persistente}
The multiset \(p_{\rm ret}(\mathcal F_\pi)\) has a unique element of
maximal multiplicity:
\[
b_\pi=(169,272,272),
\qquad
\operatorname{mult}_{\mathcal F_\pi}(b_\pi)=4.
\]
The only remaining terminal block is
\[
b_\pi^-=(418,674,521),
\]
and its multiplicity is one.
\end{proposition}

\begin{proof}
The assertion follows by exact enumeration of the five catalogue rows
satisfying \(\Psi(U)=010211\). Four rows have terminal block
\((169,272,272)\); the fifth, which carries the mutated return, has terminal
block \((418,674,521)\). The maximal multiplicity is four and is attained
by exactly one block.
\end{proof}

The oriented difference between the mutated and persistent returns is
\[
 \omega_\pi=b_\pi^- - b_\pi=(249,402,249),
 \qquad
 \langle(1,1,1),\omega_\pi\rangle=900.
\]
This identity distinguishes two objects of different types. The integer
\(169\) is a coordinate of a terminal datum and therefore belongs to degree
zero. The vector \(\omega_\pi\) admits a reading as a degree-one cochain
once the edge joining the two blocks has been oriented. Its classification
as a cocycle additionally requires a specified cochain complex and
verification that its coboundary vanishes. The analytic prolongation
constructed below uses the terminal datum independently of that additional
identification.

The selection of \(169\) is determined by the coordinate multiplicities
of the block. The symmetric group \(\mathfrak S_3\) acts by permuting its
three coordinates. The stabilizer of \(b_\pi\) exchanges the two occurrences
of \(272\) and fixes the unique coordinate index whose value has
multiplicity one. For a block of type \((r,s,s)\), with \(r\ne s\), denote
that index by \(\sigma(b)\), and define its selected value by
\[
\operatorname{sing}(b)=b_{\sigma(b)}.
\]
For \(\tau\in\mathfrak S_3\), the coordinate-index selector is equivariant,
whereas its selected value is invariant:
\[
\sigma(\tau\cdot b)=\tau\bigl(\sigma(b)\bigr),
\qquad
\operatorname{sing}(\tau\cdot b)=\operatorname{sing}(b).
\]
Equivalently, the selected value is
\[
\operatorname{sing}(b)
=\sum_{j=1}^3b_j-2\operatorname{mode}(b).
\]
In particular,
\[
\boxed{\operatorname{sing}(b_\pi)=169.}
\]

\begin{definition}[Residual selector of the microfibre]
\label{p12-83-c13-s3:def:selector-residual-pi}
The selector \(\rho_\pi\) applies two successive operations:
\begin{enumerate}
  \item it selects the terminal block of maximal multiplicity in
  \(p_{\rm ret}(\mathcal F_\pi)\);
  \item it selects the coordinate index fixed by the stabilizer of that
  block and retains the value at that index.
\end{enumerate}
By \cref{p12-83-c13-s3:prop:bloque-terminal-persistente},
\[
\boxed{\rho_\pi=169.}
\]
\end{definition}

\begin{theorem}[Equivariant uniqueness of the return selector]
\label{p12-83-c13-s3:thm:unicidad-retorno-169}
Among selectors that
\begin{enumerate}
  \item are invariant under permutations of the five regions;
  \item select a coordinate index equivariantly under permutations of the
  three terminal coordinates;
  \item select the block of maximal multiplicity;
  \item select a coordinate index fixed by the stabilizer of that block,
\end{enumerate}
the return value is uniquely determined and equals \(169\).
\end{theorem}

\begin{proof}
Invariance under \(\mathfrak S_5\) makes the selection independent of row
ordering. The third condition selects \(b_\pi\), the unique modal block
by \cref{p12-83-c13-s3:prop:bloque-terminal-persistente}. Its stabilizer in
\(\mathfrak S_3\) has two orbits on the coordinate indices: the pair of
indices carrying \(272\), and the singleton carrying \(169\). The fourth
condition selects the singleton. Equivariance transports its index under
a permutation of coordinates; the value carried by that index remains
\(169\).
\end{proof}

This argument refines the positional reading of the return. The integer
\(169\) is an invariant of the regional multiset and its reordering
symmetries. Its occurrence in the first coordinate of the displayed
persistent block is one representation of that invariant selection.
Its global frequency supplies a separate statistic: across the \(468\)
catalogue emissions, \(169\) occurs \(84\) times, distributed uniformly
with \(14\) occurrences in each of the six positions. The uniqueness in
\cref{p12-83-c13-s3:thm:unicidad-retorno-169} is therefore relational and
regional: it is determined by the combination of the microfibre, modal
persistence, and the stabilizer, independently of the isolated frequency
of the integer.

\subsubsection{Cylindrical length and formal analytic degree}

Let \(\mathcal P\) be the free module generated by regional words, with
length filtration \(F^n\mathcal P\). Its associated grading records the
first level at which each word appears. For a formal parameter \(z\), the
ordinary length character is
\[
\chi_z([w])=z^{|w|}.
\]
Its multiplicativity,
\[
\chi_z([uv])=\chi_z([u])\chi_z([v]),
\]
expresses the additivity of length under concatenation. The formal degree
of a word of length \(n\) is \(n\), represented by the monomial \(z^n\).
Subsequent evaluation at \(z=\alpha\) assigns that monomial the numerical
weight \(\alpha^n\).

\begin{definition}[Degree compatibility]
\label{p12-83-c13-s3:def:compatibilidad-grado}
The regional analytic reader is compatible with the cylindrical
filtration when it sends the homogeneous component of length \(n\) to
the formal homogeneous component of degree \(n\):
\[
\operatorname{gr}_n\mathcal P\longrightarrow
\mathbb R z^n.
\]
Evaluation at \(z=\alpha\) is a subsequent operation on the formal
expression.
\end{definition}

For every word \(w\in\mathcal F_\pi\subset\mathcal U_6\), the common
return datum enters at degree six:
\[
\rho_\pi\chi_z([w])=169z^6
\xrightarrow{\;z=\alpha\;}169\alpha^6.
\]
This term records the common impulse of the microfibre exactly once.
This six is the length of the regional word. The closed walks of
\cref{regcl27:seccion} have a different length: they jointly traverse the
four anchors
\(U_\pi,U_e,U_\varphi,U_{\rm APP}\), and their minimum length is eight.
A word, an edge of the block graph, and a closed walk are objects of
different types; the grading used here belongs to the word.

\subsubsection{The determinantal character of the common transport}

Dodecaphasic closure provides \(\alpha\). Completion of the three nonadic
pairs provides the scale \(10^3\), so that
\[
A=10^3\alpha.
\]
Once the Archimedean angular chart has been chosen, the common coordinate
in radians is
\[
x=\frac{\pi A}{180}=\frac{50\pi}{9}\alpha.
\]
Let \(R\) be a real trace-zero involution in dimension two,
\[
R^2=I_2,
\qquad
\operatorname{tr}R=0,
\]
and consider the formal two-sheet transport
\[
\mathcal T(x,y)=\exp(-xI_2-yR).
\]
The oriented variable \(y\) remains unevaluated. The action of
\(\mathcal T\) on the determinant line is nevertheless independent of
that variable:
\[
\begin{aligned}
\det\mathcal T(x,y)
&=\exp\!\left(\operatorname{tr}(-xI_2-yR)\right)\\
&=e^{-2x}.
\end{aligned}
\]
Define
\[
\boxed{
D_A=e^{-2x}=e^{-100\pi\alpha/9}.}
\]
The subscript identifies the determinant of the common contraction
associated with \(A\). Its definition depends on that common coordinate,
while \(C_*^{\rm HMT}\) belongs to the conjugate angular reading.

\begin{proposition}[Character of the determinant line]
\label{p12-83-c13-s3:prop:caracter-determinante}
The map
\[
\delta_A:(\mathbb N_0,+)\longrightarrow(\mathbb R_{>0},\cdot),
\qquad
\delta_A(k)=D_A^k,
\]
is the unique multiplicative character whose value on an elementary
prolongation is \(D_A\).
\end{proposition}

\begin{proof}
Multiplicativity gives
\[
\delta_A(k)
=\delta_A(\underbrace{1+\cdots+1}_{k\text{ terms}})
=\delta_A(1)^k=D_A^k.
\]
The same identity establishes existence and uniqueness.
\end{proof}

The determinant is invariant under conjugation of \(\mathcal T\) and
under the exchange of sheets \(R\mapsto-R\). Accordingly, the tail
constructed from \(D_A\) belongs to the common mode. Orientation enters
through the sign with which the regional character corrects the phase.

\Needspace{29\baselineskip}
\subsubsection{Rules of the regional analytic reader}

The analytic extension is specified by the following rules.

\begin{definition}[Determinantal regional reader]
\label{p12-83-c13-s3:def:lector-regional-determinantal}
The determinantal regional reader satisfies:
\begin{enumerate}
  \item \emph{degree compatibility}: cylindrical length determines the
  formal analytic degree;
  \item \emph{renewal decomposition}: the finite return impulse and the
  strict continuation belong to distinct additive summands;
  \item \emph{continuation normalization}: after the residue
  \(\rho_\pi\) has been recorded once, the surviving branch begins at the
  unit of the determinant line;
  \item \emph{concatenation covariance}: each elementary prolongation
  acts through
  \(\bigwedge^2\mathcal T=D_A\,\operatorname{Id}_{\bigwedge^2\mathbb R^2}\);
  \item \emph{orientation}: in the fixed cardinal convention for APP,
  the regional return belongs to the compensating sector and is
  subtracted from the fifth-degree action phase.
\end{enumerate}
\end{definition}

The third rule has mathematical content. The value \(169\) is recorded
once as the return amplitude. Subsequent prolongations propagate the
strict continuation, which counts the unique continuing line and
normalizes it by its unit. This distinguishes the renewal recurrence
from homogeneous iteration of the initial impulse.

The need to specify this normalization can be demonstrated. For every
\(\lambda\in\mathbb R\), the collection
\[
F_\lambda(z)
=169z^6+\lambda\frac{D_Az^7}{1-D_Az}
\]
preserves the same finite return and the same multiplicative recurrence for
the tail coefficients. The catalogue and that recurrence leave
\(\lambda\) undetermined. The unit of the determinant line fixes
\(\lambda=1\) before evaluation at \(z=\alpha\). Thus normalization is
part of the definition of the character and precedes its final numerical
evaluation.
